\documentclass[11pt,a4paper]{article}

% ── Codificación, fuentes y lenguaje ───────────────────────────────────────────
\usepackage[utf8]{inputenc}
\usepackage[T1]{fontenc}
\usepackage[default,scale=0.95]{sourcesanspro}
\renewcommand{\familydefault}{\sfdefault}
\usepackage[spanish,es-noshorthands,es-tabla]{babel}

% ── Geometría y espaciado ─────────────────────────────────────────────────────
\usepackage[top=2.2cm, bottom=2.5cm, left=2.2cm, right=2.2cm, headheight=15pt]{geometry}
\usepackage{setspace}
\setstretch{1.18}
\usepackage{parskip}

% ── Matemáticas y unidades ────────────────────────────────────────────────────
\usepackage{amsmath}
\usepackage{amssymb}
\usepackage{siunitx}

% ── Circuitos ─────────────────────────────────────────────────────────────────
\usepackage[european, straightvoltages]{circuitikz}

% ── Tablas ────────────────────────────────────────────────────────────────────
\usepackage{booktabs}
\usepackage{tabularx}
\usepackage{array}
\usepackage{longtable}
\usepackage{colortbl}
\usepackage{enumitem}

% ── Colores, cajas e iconos ───────────────────────────────────────────────────
\usepackage[dvipsnames]{xcolor}
\usepackage{tcolorbox}
\tcbuselibrary{skins, breakable}
\usepackage{fontawesome5}
\usepackage{graphicx}

% ── Listados de código (dark-mode infográfico) ────────────────────────────────
%   Estilo base: fondo oscuro con números de línea. Los bloques lstlisting se
%   envuelven automáticamente en una caja tcolorbox redondeada.
\usepackage{listings}

% ── Secciones con barra de color (infografía) ────────────────────────────────
\usepackage{titlesec}
\titleformat{\section}
  {\normalfont\LARGE\bfseries\color{azulNoche}}
  {\colorbox{cyanNeon}{\color{fondoOscuro}\thesection}}{0.7em}{}
  [\vspace{2pt}\color{cyanNeon}\rule{\linewidth}{1.8pt}]
\titlespacing*{\section}{0pt}{24pt}{10pt}
\titleformat{\subsection}
  {\normalfont\large\bfseries\color{azulNoche}}
  {\textcolor{cyanNeon}{\thesubsection}}{0.7em}{}
  [\vspace{1pt}\color{grisLinea}\rule{\linewidth}{0.6pt}]
\titlespacing*{\subsection}{0pt}{14pt}{6pt}
\titleformat{\subsubsection}
  {\normalfont\normalsize\bfseries\color{azulNoche}}
  {\textcolor{cyanNeon}{\thesubsubsection}}{0.7em}{}
\titlespacing*{\subsubsection}{0pt}{10pt}{4pt}

% ── Cabeceras y pies de página ────────────────────────────────────────────────
\usepackage{fancyhdr}
\usepackage{lastpage}
\pagestyle{fancy}
\fancyhf{}
\renewcommand{\headrulewidth}{0pt}
\renewcommand{\footrulewidth}{0pt}
\fancyfoot[C]{%
  \begin{minipage}{\linewidth}
    \vspace{2pt}
    \color{cyanNeon}\rule{\linewidth}{1.2pt}\\[2pt]
    \centering\color{grisTexto}\footnotesize
    Página \thepage\ de \pageref{LastPage}
  \end{minipage}%
}

% ── Hiperenlaces ──────────────────────────────────────────────────────────────
\usepackage[colorlinks=true, linkcolor=azulNoche, urlcolor=cyanNeon]{hyperref}
\sloppy
\emergencystretch 3em

% ── Paleta de colores — tecnológico dark-mode ──────────────────────────────────
\definecolor{fondoOscuro}{HTML}{0D1B2A}     % Dark navy — fondo banda título
\definecolor{verdeTurquesa}{HTML}{004D40}    % Verde turquesa — bandas de marca
\definecolor{azulNoche}{HTML}{1B3A6B}       % Midnight blue — secciones, marcos
\definecolor{azulMedio}{HTML}{2A5298}       % Azul medio — variante de acento
\definecolor{cyanNeon}{HTML}{00C8E0}        % Cyan eléctrico — acentos primarios
\definecolor{verdeSignal}{HTML}{00B887}     % Verde señal — fortalezas, éxito
\definecolor{naranjaVivo}{HTML}{FF7043}     % Naranja vivo — mejoras, advertencias
\definecolor{rojoAlerta}{HTML}{E53935}      % Rojo alerta — errores
\definecolor{grisPapel}{HTML}{F4F6FA}       % Fondo tarjetas claras
\definecolor{grisLinea}{HTML}{C8D0E0}       % Bordes sutiles
\definecolor{grisTexto}{HTML}{6B7A99}       % Texto secundario
\definecolor{amarilloNota}{HTML}{FFB300}    % Notas / advertencias doradas

% Colores específicos para bloques de código (dark-mode)
\definecolor{codigoFondo}{HTML}{1E2D3D}      % Fondo del bloque de código
\definecolor{codigoComentario}{HTML}{7A8FA6} % Comentarios
\definecolor{codigoCadena}{HTML}{FF9D5C}     % Cadenas / strings
\definecolor{codigoKeyword}{HTML}{00C8E0}    % Palabras clave
\definecolor{codigoNumero}{HTML}{00E5A0}     % Números

% ── Banda de título: portada infográfica ──────────────────────────────────────
%   Diseño en 2 capas: banda de color (por defecto fondo oscuro) + franja de
%   acento cyan abajo. Uso: \bandaTitulo[color]{título}{subtítulo}.
%   El icono se puede cambiar con \renewcommand{\iconoBanda}{...} (FontAwesome).
\newcommand{\iconoBanda}{microchip}
\newcommand{\bandaTitulo}[3][fondoOscuro]{%
  \begin{tcolorbox}[
    enhanced,
    colback=#1, colframe=#1,
    boxrule=0pt, arc=8pt, outer arc=8pt,
    width=\linewidth,
    top=18pt, bottom=6pt, left=14pt, right=14pt,
    borderline south={3.5pt}{0pt}{cyanNeon},
  ]
    \begin{minipage}[c]{0.07\linewidth}
      \centering
      \textcolor{cyanNeon}{\fontsize{30}{30}\selectfont\faIcon{\iconoBanda}}
    \end{minipage}%
    \hspace{8pt}%
    \begin{minipage}[c]{0.88\linewidth}
      {\fontsize{20}{24}\selectfont\bfseries\color{white}#2}\par\vspace{5pt}%
      \textcolor{cyanNeon!80!white}{\small\faIcon{angle-right}~#3}%
    \end{minipage}
    \vspace{4pt}
  \end{tcolorbox}%
  \vspace{8pt}%
}

% ── Tarjetas de datos (infografía) ────────────────────────────────────────────
\tcbset{
  tarjetaDato/.style={
    enhanced, breakable,
    arc=6pt, outer arc=6pt,
    colback=grisPapel, colframe=azulNoche,
    boxrule=1pt,
    fonttitle=\bfseries\small, coltitle=white,
    attach boxed title to top left={yshift=-3mm, xshift=6mm},
    boxed title style={
      colback=azulNoche, colframe=azulNoche,
      arc=4pt, boxrule=0pt,
    },
    drop fuzzy shadow=azulNoche!30!white,
    top=8pt, bottom=8pt, left=10pt, right=10pt,
  },
  tarjetaIcono/.style={
    enhanced, breakable,
    arc=6pt, outer arc=6pt,
    colback=white, colframe=grisLinea, boxrule=0.8pt,
    fonttitle=\bfseries\small, coltitle=azulNoche,
    borderline west={3pt}{0pt}{cyanNeon},
    drop fuzzy shadow=azulNoche!25!white,
    top=6pt, bottom=6pt, left=10pt, right=10pt,
  },
  cajaTitulo/.style={
    enhanced, breakable,
    arc=6pt, outer arc=6pt,
    colback=azulNoche!10!grisPapel, colframe=azulNoche,
    boxrule=1pt,
    fonttitle=\bfseries, coltitle=white,
    attach boxed title to top left={yshift=-3mm, xshift=6mm},
    boxed title style={
      colback=azulNoche, colframe=azulNoche,
      arc=4pt, boxrule=0pt,
    },
    drop fuzzy shadow=azulNoche!25!white,
    top=8pt, bottom=8pt, left=10pt, right=10pt,
  },
  cajaContenido/.style={
    enhanced, breakable,
    arc=5pt, outer arc=5pt,
    colback=grisPapel, colframe=grisLinea,
    boxrule=0.8pt,
    fonttitle=\bfseries\small, coltitle=azulNoche,
    top=6pt, bottom=6pt, left=10pt, right=10pt,
  },
  cajaFortaleza/.style={
    enhanced, breakable,
    arc=6pt, outer arc=6pt,
    colback=verdeSignal!8!white, colframe=verdeSignal,
    boxrule=1pt,
    borderline west={4pt}{0pt}{verdeSignal},
    fonttitle=\bfseries\small, coltitle=verdeSignal!70!black,
    drop fuzzy shadow=verdeSignal!20!white,
    top=6pt, bottom=6pt, left=10pt, right=10pt,
  },
  cajaMejora/.style={
    enhanced, breakable,
    arc=6pt, outer arc=6pt,
    colback=naranjaVivo!8!white, colframe=naranjaVivo,
    boxrule=1pt,
    borderline west={4pt}{0pt}{naranjaVivo},
    fonttitle=\bfseries\small, coltitle=naranjaVivo!80!black,
    drop fuzzy shadow=naranjaVivo!20!white,
    top=6pt, bottom=6pt, left=10pt, right=10pt,
  },
  cajaRecomendacion/.style={
    enhanced, breakable,
    arc=6pt, outer arc=6pt,
    colback=cyanNeon!6!white, colframe=cyanNeon!60!azulNoche,
    boxrule=1pt,
    borderline west={4pt}{0pt}{cyanNeon},
    fonttitle=\bfseries\small, coltitle=azulNoche,
    drop fuzzy shadow=azulNoche!20!white,
    top=6pt, bottom=6pt, left=10pt, right=10pt,
  },
}

% ── Ícono + texto (encabezados de sección y elementos) ────────────────────────
\newcommand{\iconotexto}[2]{\textcolor{cyanNeon}{\faIcon{#1}}~#2}

% ── Listados de código: estilo dark + auto-envoltura ─────────────────────────
\lstdefinestyle{estiloCodigo}{
    backgroundcolor=\color{codigoFondo},
    basicstyle=\ttfamily\small\color{white},
    commentstyle=\color{codigoComentario}\itshape,
    stringstyle=\color{codigoCadena},
    keywordstyle=\color{codigoKeyword}\bfseries,
    numberstyle=\tiny\color{codigoComentario},
    numbers=left,
    numbersep=10pt,
    showstringspaces=false,
    breaklines=true,
    breakatwhitespace=false,
    tabsize=4,
    frame=none,
    captionpos=b,
    aboveskip=6pt,
    belowskip=4pt,
    xleftmargin=14pt,
    framexleftmargin=14pt,
    literate=
      {á}{{\'a}}1 {é}{{\'e}}1 {í}{{\'i}}1 {ó}{{\'o}}1 {ú}{{\'u}}1
      {Á}{{\'A}}1 {É}{{\'E}}1 {Í}{{\'I}}1 {Ó}{{\'O}}1 {Ú}{{\'U}}1
      {ñ}{{\~n}}1 {Ñ}{{\~N}}1 {ü}{{\"u}}1 {Ü}{{\"U}}1
      {¿}{{?`}}1 {¡}{{!`}}1
}
\lstdefinestyle{estiloPython}{
    style=estiloCodigo,
    language=Python,
}
\lstset{style=estiloCodigo}
\BeforeBeginEnvironment{lstlisting}{%
  \begin{tcolorbox}[
    enhanced, breakable,
    arc=6pt, outer arc=6pt,
    colback=codigoFondo, colframe=azulNoche,
    boxrule=1pt,
    drop fuzzy shadow=azulNoche!25!white,
    top=2pt, bottom=2pt, left=0pt, right=0pt,
  ]%
}
\AfterEndEnvironment{lstlisting}{\end{tcolorbox}}

% ── Cajas de contenido temáticas ──────────────────────────────────────────────
\newtcolorbox{cajaRecuerda}{
    enhanced, breakable,
    arc=6pt, outer arc=6pt,
    colback=azulNoche!10!grisPapel,
    colframe=azulNoche, boxrule=1pt,
    borderline west={4pt}{0pt}{cyanNeon},
    fonttitle=\bfseries\small\color{white},
    title={\faIcon{info-circle}~Recuerda},
    attach boxed title to top left={yshift=-3mm, xshift=6mm},
    boxed title style={colback=azulNoche, arc=4pt, boxrule=0pt},
    drop fuzzy shadow=azulNoche!25!white,
    left=10pt, right=10pt, top=8pt, bottom=8pt,
}

\newtcolorbox{cajaConcepto}{
    enhanced, breakable,
    arc=6pt, outer arc=6pt,
    colback=verdeSignal!8!white,
    colframe=verdeSignal, boxrule=1pt,
    borderline west={4pt}{0pt}{verdeSignal},
    fonttitle=\bfseries\small\color{white},
    title={\faIcon{lightbulb}~Concepto clave},
    attach boxed title to top left={yshift=-3mm, xshift=6mm},
    boxed title style={colback=verdeSignal!80!black, arc=4pt, boxrule=0pt},
    drop fuzzy shadow=verdeSignal!20!white,
    left=10pt, right=10pt, top=8pt, bottom=8pt,
}

\newtcolorbox{cajaEjemplo}{
    enhanced, breakable,
    arc=6pt, outer arc=6pt,
    colback=cyanNeon!5!white,
    colframe=cyanNeon!70!azulNoche, boxrule=1pt,
    borderline west={4pt}{0pt}{cyanNeon},
    fonttitle=\bfseries\small\color{fondoOscuro},
    title={\faIcon{code}~Ejemplo},
    attach boxed title to top left={yshift=-3mm, xshift=6mm},
    boxed title style={colback=cyanNeon, arc=4pt, boxrule=0pt},
    drop fuzzy shadow=azulNoche!20!white,
    left=10pt, right=10pt, top=8pt, bottom=8pt,
}

\newtcolorbox{cajaEjercicio}{
    enhanced, breakable,
    arc=6pt, outer arc=6pt,
    colback=azulMedio!7!white,
    colframe=azulMedio, boxrule=1pt,
    borderline west={4pt}{0pt}{azulMedio},
    fonttitle=\bfseries\small\color{white},
    title={\faIcon{pencil-alt}~Ejercicio},
    attach boxed title to top left={yshift=-3mm, xshift=6mm},
    boxed title style={colback=azulMedio, arc=4pt, boxrule=0pt},
    drop fuzzy shadow=azulNoche!20!white,
    left=10pt, right=10pt, top=8pt, bottom=8pt,
}

\newtcolorbox{cajaReto}{
    enhanced, breakable,
    arc=6pt, outer arc=6pt,
    colback=naranjaVivo!6!white,
    colframe=naranjaVivo, boxrule=1pt,
    borderline west={4pt}{0pt}{naranjaVivo},
    fonttitle=\bfseries\small\color{white},
    title={\faIcon{fire}~Reto},
    attach boxed title to top left={yshift=-3mm, xshift=6mm},
    boxed title style={colback=naranjaVivo, arc=4pt, boxrule=0pt},
    drop fuzzy shadow=naranjaVivo!20!white,
    left=10pt, right=10pt, top=8pt, bottom=8pt,
}

% ── Indicadores de dificultad ─────────────────────────────────────────────────
\newcommand{\facil}{\textcolor{verdeSignal}{\faIcon{circle}\,\textbf{Fácil}}}
\newcommand{\intermedio}{\textcolor{naranjaVivo}{\faIcon{circle}\,\textbf{Intermedio}}}
\newcommand{\dificil}{\textcolor{rojoAlerta}{\faIcon{circle}\,\textbf{Difícil}}}

% ─────────────────────────────────────────────────────────────────────────────

% ── Cabeceras específicas del documento ────────────────────────────────────────
\fancyhead[L]{\small\color{azulNoche}\textbf{ELECTRÓNICA} --- Higiene de Estado de Sesión}
\fancyhead[R]{\small\color{grisTexto}11/09/2026}
\renewcommand{\headrulewidth}{0pt}

% ─────────────────────────────────────────────────────────────────────────────
\begin{document}

% ══ PORTADA INFográfica ════════════════════════════════════════════════════════
\renewcommand{\iconoBanda}{shield-alt}
\bandaTitulo[fondoOscuro]{Higiene de Estado de Sesión}{run\_state.json como vista derivada, guards en los servidores MCP, limpieza determinista y eslabón semántico de bitácoras}

\vspace{4pt}
\begin{tcolorbox}[cajaTitulo, title=Resumen ejecutivo]
\textbf{Problema:} \texttt{run\_state*.json} son \emph{vistas derivadas} de los logs
append-only \texttt{session\_log\_*.jsonl} (la fuente de verdad). Una vista puede
\emph{sobrevivir a su corrida} y quedar huérfana, contaminando las respuestas de los
servidores MCP con el estado de \emph{otra} ejecución. Además, la memoria semántica
del proyecto (por QUÉ se decidió, qué se descartó) vive \emph{exclusivamente} en
las bitácoras de sesión \texttt{Sessions/}, y no había mecanismo determinista para
asegurar que existieran ni que tuvieran la estructura mínima.

\textbf{Solución (aprobada por el usuario):} dos pillars de higiene, ambos a
\textbf{0 créditos} OpenRouter:
\begin{enumerate}[leftmargin=*]
  \item \textbf{Bajo nivel (estado de ejecución):} script determinista
        (\texttt{execution/estado\_sesion.py}) con \texttt{check} + \texttt{clean},
        una guardia por mtime en los 5 servidores MCP, edge case documentado en la
        directiva, y regla de \emph{state hygiene} en \texttt{AGENTS.md}.
  \item \textbf{Alto nivel (contexto semántico):} script determinista
        (\texttt{execution/bitacoras.py}) con \texttt{nueva --tema} (plantilla
        canónica que nunca sobrescribe) y \texttt{check [--today]} (valida nombre,
        fecha y secciones mínimas de las bitácoras recientes), edge case de ``bitácora
        no escrita'' en la directiva, y regla de cierre en \texttt{AGENTS.md}.
\end{enumerate}
\end{tcolorbox}

% ══ SECCIÓN 1: EL PROBLEMA ═════════════════════════════════════════════════════
\section{\iconotexto{bullseye}{El problema: vistas que sobreviven al run}}

\subsection{Modelo de estado en ELECTRONICA}
El rastreo de flujos multi-paso sigue una separación estricta entre \emph{fuente de
verdad} y \emph{vista}:

\begin{center}
\begin{tabularx}{\linewidth}{>{\raggedright\arraybackslash}p{0.30\linewidth} p{0.30\linewidth} X}
\toprule
\textbf{Artefacto} & \textbf{Rol} & \textbf{Naturaleza} \\
\midrule
\texttt{session\_log\_<run>.jsonl} & Fuente de verdad & Append-only, inmutable \\
\texttt{run\_state.json} & Vista global del último flujo & Sobrescribible \\
\texttt{run\_state\_<run>.json} & Vista de un run concreto & Regenerable \\
\bottomrule
\end{tabularx}
\end{center}

\begin{cajaConcepto}
\faIcon{lightbulb}~\textbf{Concepto clave:} un \texttt{run\_state.json} puede
\emph{quedar huérfano}. Si el flujo termina (evento \texttt{flujo/fin} en su log) pero
la vista no se limpia, el archivo sigue existiendo en \texttt{.tmp/} apuntando a un run
que ya no está activo. Cualquier consumidor que lo lea \emph{sin verificar} reportará
estado de una corrida terminada — o peor, de una corrida distinta.
\end{cajaConcepto}

\subsection{Caso real detectado}
Durante la auditoría de continuidad se encontró exactamente este caso:

\begin{tcolorbox}[cajaMejora, title=Caso de estudio --- vista huérfana de un dry-run]
\begin{itemize}%[leftmargin=*]
  \item \textbf{\texttt{.tmp/run\_state.json}} apuntaba a
        \texttt{flujo-repo-a-skill-2026-09-09T12-29-54} (un dry-run de prueba de
        trazabilidad).
  \item Su log \texttt{session\_log\_...jsonl} tenía \textbf{6 eventos} y terminaba con
        \texttt{flujo/fin}: la corrida \emph{ya había terminado}.
  \item El estado referenciaba \texttt{deepseek/deepseek-v4-pro}, modelo
        \textbf{discontinuado} (sustituido por \texttt{deepseek-v4.1-flash}).
  \item Los servidores MCP leen \texttt{run\_state.json} \emph{global} como
        ``estado del flujo recién lanzado'' → riesgo real de responder con el contexto
        de otra corrida.
\end{itemize}
\end{tcolorbox}

La trazabilidad append-only estaba íntegra: la cadena de hashes del log era válida
(\texttt{integrity} OK). El fallo no era del log, sino de \emph{confiar en la vista sin
cruzar contra la fuente}.

% ══ SECCIÓN 2: LA SOLUCIÓN ═════════════════════════════════════════════════════
\section{\iconotexto{layer-group}{La solución en tres capas}}

Siguiendo la arquitectura determinista, el refuerzo cubrió las tres capas:

\begin{table}[h]
\centering
\begin{tabularx}{\linewidth}{>{\raggedright\arraybackslash}p{0.20\linewidth} p{0.36\linewidth} X}
\toprule
\textbf{Capa} & \textbf{Artefacto} & \textbf{Aporte} \\
\midrule
L3 Ejecución & \texttt{execution/estado\_sesion.py} & Diagnóstico y purga de vistas huérfanas (\texttt{check}/\texttt{clean}) \\
L3 Ejecución & \texttt{execution/bitacoras.py} & Plantilla (\texttt{nueva}) y validación de bitácoras semánticas (\texttt{check}) \\
L1 Directiva & \texttt{directives/trazabilidad\_sesiones.yaml} & Edge cases ``run\_state huérfano'' y ``bitácora sin escribir'' \\
L2 Orquestación & \texttt{AGENTS.md} + guards en MCP & Higiene de bajo nivel al arrancar; cierre de bitácora al finalizar \\
\bottomrule
\end{tabularx}
\end{table}

\subsection{Capa de ejecución: execution/estado\_sesion.py}

\subsubsection{check --- diagnóstico, nunca borra}
Escanea \texttt{run\_state.json} y \texttt{run\_state\_*.json}, cruza cada vista contra
su log y emite un veredicto por artefacto:

\begin{lstlisting}[style=estiloPython]
$ python3 execution/estado_sesion.py check
{
  "checks": [ { "archivo": "..../.tmp/run_state.json",
                "run_id": "flujo-repo-a-skill-...",
                "veredicto": "huerfano",
                "razon": "la corrida termino (flujo/fin); la vista sobrevivio.",
                "modelo_obsoleto": "deepseek-v4-pro" } ],
  "veredicto_global": "atencion", "anomalias": 1
}
\end{lstlisting}

Veredictos posibles y significado:

\begin{center}
\begin{tabularx}{\linewidth}{>{\raggedright\arraybackslash}p{0.14\linewidth} X}
\toprule
\textbf{Veredicto} & \textbf{Significado} \\
\midrule
\texttt{huerfano} & El log existe y su último evento es \texttt{flujo/fin}: la vista sobrevivió a la corrida terminada \\
\texttt{vigente} & Log sin \texttt{flujo/fin}: corrida en curso o reanudable \\
\texttt{no\_verificable} & No hay log (flujo sin trazabilidad): \textbf{no} se borra por falta de certeza \\
\texttt{corrupto} & JSON inválido o sin \texttt{run\_id} (escritura a medias) \\
\bottomrule
\end{tabularx}
\end{center}

Además detecta \textbf{modelos obsoletos} por substring (\texttt{deepseek-v4-pro}) y lo
reporta como \texttt{advertencia} (no justifica borrado por sí solo).

\subsubsection{clean --- borra solo con certeza}
\begin{lstlisting}[style=estiloPython]
$ python3 execution/estado_sesion.py clean            # purga vistas huerfanas
$ python3 execution/estado_sesion.py clean --dry-run  # lista sin borrar
$ python3 execution/estado_sesion.py check --run <id> # acotar a un run
\end{lstlisting}

Regla de oro del \texttt{--clean}: \textbf{borra únicamente} las vistas cuya corrida
terminó con \texttt{flujo/fin} (certeza absoluta). \textbf{Nunca} borra los logs
append-only (fuente de verdad inmutable), y \textbf{no} toca los
\texttt{no\_verificable}.

\begin{cajaRecuerda}
\faIcon{info-circle}~\textbf{Recuerda:} \texttt{check} solo diagnostica. El borrado es
intencional, vía \texttt{clean}, y está restringido a lo que el criterio anterior puede
confirmar. Un flujo sin trazabilidad (\texttt{no\_verificable}) se advierte pero jamás
se purga automáticamente.
\end{cajaRecuerda}

\subsection{Capa de directiva: trazabilidad\_sesiones.yaml}
Se amplió la directiva con dos edge cases específicos (protocolos de recuperación):

\begin{lstlisting}
- case: "run_state.json obsoleto o huerfano al inicio de una sesion/consulta MCP"
  recovery: >
    Ejecutar execution/estado_sesion.py check y purgar con clean si marca
    huerfano/corrupto. Nunca tocar los logs (fuente inmutable).

- case: "No se escribio la bitacora semantica de la sesion"
  recovery: >
    Al iniciar: si no existe bitacora del dia, crearla con
      python3 execution/bitacoras.py nueva --tema "<tema>".
    Al cerrar: rellenarla y validar con bitacoras.py check.
    El contexto de ALTO nivel (por que se decidio, que se descarto) vive
    EXCLUSIVAMENTE en Sessions/. Sin bitacora, la proxima sesion pierde
    contexto y volveria a investigar desde cero.
\end{lstlisting}

\texttt{bitacoras.py} está referenciado en la lista de scripts de la directiva
y el edge case cubre la continuidad semántica entre sesiones.

\subsection{Capa de orquestación: guards en los MCP y AGENTS.md}

\subsubsection{Guardia por mtime en los 5 servidores MCP}
Los servidores \texttt{mcp\_evaluar}, \texttt{mcp\_diagnostico}, \texttt{mcp\_analizar},
\texttt{mcp\_elaborar} y \texttt{mcp\_docs} siguen el mismo patrón: lanzan el flujo con
\texttt{subprocess.run} y, si termina OK, leen \texttt{run\_state.json} para dar una
respuesta \emph{rica}. La corrección captura el mtime \emph{antes} de lanzar y solo lee
la vista si el archivo fue \textbf{reescrito durante esta corrida}:

\begin{lstlisting}[style=estiloPython]
state_file = os.path.join(project_root, ".tmp", "run_state.json")
mtime_before = (os.path.getmtime(state_file)
                if os.path.exists(state_file) else None)
try:
    resultado = subprocess.run(cmd, capture_output=True, text=True,
                               encoding="utf-8", cwd=project_root)
    state_data = {}
    if os.path.exists(state_file):
        try:
            if mtime_before is not None and \
               os.path.getmtime(state_file) <= mtime_before:
                pass  # vista de OTRA corrida: no reportar
            else:
                with open(state_file, "r", encoding="utf-8") as f:
                    state_data = json.load(f)
        except Exception:
            state_data = {}
\end{lstlisting}

\begin{cajaRecuerda}
\faIcon{exclamation-triangle}~\textbf{Corrección de diseño importante:} la primera
versión del guard comprobaba \texttt{"flujo/fin"} en el log para descartar la vista.
\textbf{Ese criterio era inválido}: un flujo que corre \emph{bien} también emite
\texttt{flujo/fin}, así que se habría descartado el estado legítimo de la corrida que
acaba de terminar con éxito. El criterio correcto para el MCP es temporal:
``¿el \texttt{run\_state.json} fue reescrito \emph{durante} este lanzamiento?''
→ comparación de mtime. Se corrigió en los 5 servers antes del commit.
\end{cajaRecuerda}

\subsubsection{Regla de arranque y cierre en AGENTS.md}
\begin{lstlisting}
- **State hygiene (freshness):** run_state*.json are *derived views* of the
  append-only session_log_*.jsonl; a view can outlive its run (orphan) and poison
  MCP responses. At the start of each session, run `python3 execution/estado_sesion.py
  check`; purge confirmed orphans with `... clean` (never touches the immutable logs).

- **Session logs (continuity):** ...
  - At the start: create bitácora via `python3 execution/bitacoras.py nueva --tema "<tema>"`
  - At session close: run `... check` and fix any missing day-log or empty sections
\end{lstlisting}

\begin{cajaConcepto}
\faIcon{lightbulb}~\textbf{Dos niveles de memoria, dos instrumentos:}
el \emph{bajo nivel} (datos reproducibles: log + run\_state) se gestiona con
\texttt{estado\_sesion.py}. El \emph{alto nivel} (semántica: por qué se decidió,
qué se descartó) se gestiona con \texttt{bitacoras.py}. Sin ninguno de los dos, la
próxima sesión re-investiga desde cero (tokens/créditos quemados).
\end{cajaConcepto}

% ══ SECCIÓN 5: EL ESLABÓN SEMÁNTICO — BITÁCORAS ═════════════════════════════
\section{\iconotexto{book}{Eslabón semántico: bitácoras de sesión}}

\begin{cajaConcepto}
\faIcon{lightbulb}~\textbf{El problema:} la salud de ejecución (sección anterior) blinda
el \emph{bajo nivel}: datos reproducibles, deterministas, verificables. Pero el
\emph{contexto de alto nivel} — por qué se decidió algo, qué se descartó, las
intenciones del usuario, los matices que no caben en un log JSON — vive
\textbf{exclusivamente} en las bitácoras markdown de \texttt{Sessions/}. Si no se
escriben, la siguiente sesión \emph{no puede recuperar} ese razonamiento y volvería
a investigar desde cero, quemando tokens y tomando decisiones sin contexto.
\end{cajaConcepto}

\subsection{execution/bitacoras.py}

Script determinista (0 créditos) con dos subcomandos:

\subsubsection{nueva --tema \textless{}tema\textgreater{}}
Genera la plantilla canónica de la sesión de HOY (nunca sobrescribe un archivo
existente). Formato del nombre: \texttt{YYYY-MM-DD\_\textless{}tema\textgreater{}.md}
(slug ASCII con transliteración de acentos). Plantilla:

\begin{lstlisting}
# 2026-09-11 — Tema de la sesión

## Tema
(1 frase sobre qué se abordó)

## Contexto
(Por qué se aborda, eventos detectados)

## Decisiones (usuario)
(Acuerdos explícitos del usuario — NO inventar)

## Actividades
(Pasos, scripts, salidas, veredictos)

## Pendientes
(Queda abierto, pendientes acordados de otras sesiones)
\end{lstlisting}

\subsubsection{check [--today]}
Diagnostica la estructura de las bitácoras existentes: valida nombre
(\texttt{YYYY-MM-DD\_\textless{}tema\textgreater{}.md}), fecha parseable (sin
futuros) y las cuatro secciones mínimas. Clasifica cada bitácora en tres periodos:

\begin{center}
\begin{tabularx}{\linewidth}{>{\raggedright\arraybackslash}p{0.16\linewidth} X}
\toprule
\textbf{Periodo} & \textbf{Significado} \\
\midrule
\texttt{hoy} & Sesión del día en curso — \textbf{debe} existir y cumplir plantilla \\
\texttt{reciente} & Desde la adopción de la plantilla — debe cumplir estructura \\
\texttt{legado} & Antes de la plantilla — formato heredado, informativo, no penaliza \\
\bottomrule
\end{tabularx}
\end{center}

Solo las anomalías en \texttt{hoy}/\texttt{reciente} cuentan para
\texttt{veredicto\_global} y el exit code. El check también reporta si la sesión
del día ya dejó bitácora (\texttt{sesion\_hoy.hay\_bitacora}).

\begin{lstlisting}[style=estiloPython]
$ python3 execution/bitacoras.py check
{
  "bitacoras": [...],
  "sesion_hoy": {"fecha": "2026-09-11",
                 "hay_bitacora": true,
                 "archivos": ["Sessions/2026-09-11_....md"]},
  "anomalias": 0, "accionables": 0,
  "veredicto_global": "ok"
}
\end{lstlisting}

\subsection{Hábito codificado en AGENTS.md}
La regla en \texttt{AGENTS.md} vincula el uso del script al flujo de trabajo:

\begin{itemize}[leftmargin=*]
  \item \textbf{Al inicio:} buscar bitácora previa del tema → crear la nueva con
        \texttt{bitacoras.py nueva --tema}.
  \item \textbf{Al cerrar:} rellenar las secciones y ejecutar
        \texttt{bitacoras.py check}; si reporta anomalía, corregir antes de terminar.
\end{itemize}

\begin{cajaRecuerda}
\faIcon{info-circle}~\textbf{Las bitácoras legado no se reescriben:} las sesiones
anteriores a la adopción de la plantilla (2026-09-11) usaban otro formato. Se listan
informativas para mantener trazabilidad histórica; solo las nuevas (desde hoy) deben
seguir la plantilla canónica.
\end{cajaRecuerda}

% ══ SECCIÓN 6: VERIFICACIÓN ═══════════════════════════════════════════════════
\section{\iconotexto{flask}{Verificación (0 créditos)}}

\subsection{Guardia MCP: dos casos controlados}
Se replicó la lógica exacta del server contra dos escenarios:

\begin{tcolorbox}[cajaFortaleza, title=Caso A --- vista huérfana (no reescrita)]
\begin{lstlisting}[style=estiloPython]
run_state.json escrito hace 1h (mtime viejo).
mtime_before = getmtime()  # la corrida NO reescribe el archivo
resultado: state_data == {}       # << no se reporta estado ajeno
\end{lstlisting}
Resultado: \textbf{se descarta} la vista de otra corrida. \faIcon{check-circle}
\end{tcolorbox}

\begin{tcolorbox}[cajaFortaleza, title=Caso B --- estado reescrito por la corrida]
\begin{lstlisting}[style=estiloPython]
mtime_before = getmtime()
el flujo reescribe run_state.json  # mtime > mtime_before
resultado: state_data == {"run_id": "...", "context": {...}}  # << se usa
\end{lstlisting}
Resultado: \textbf{se conserva} el estado legítimo. \faIcon{check-circle}
\end{tcolorbox}

\subsection{Limpieza real del caso detectado}
\begin{lstlisting}[style=estiloPython]
$ python3 execution/estado_sesion.py check
  # -> veredicto "huerfano" + advertencia deepseek-v4-pro
$ python3 execution/estado_sesion.py clean
  # -> eliminado: .tmp/run_state.json
$ python3 execution/estado_sesion.py check
  # -> checks: []  veredicto_global: ok   (log immutable intacto)
\end{lstlisting}

El log append-only \texttt{session\_log\_flujo-repo-a-skill-2026-09-09T12-29-54.jsonl}
permaneció \textbf{intacto} (fuente de verdad intacta).

% ══ SECCIÓN 7: ALINEACIÓN CON LA ARQUITECTURA ═════════════════════════════════
\section{\iconotexto{sitemap}{Alineación con la arquitectura}}

\begin{itemize}%[leftmargin=*]
  \item \textbf{Determinismo:} mismo estado de \texttt{.tmp/} → mismo veredicto. El
        script no razona ni consume créditos.
  \item \textbf{Separación de responsabilidades:} la decisión de \emph{cuándo} una
        vista es obsoleta o una bitácora está incompleta vive en código (capas 1 y 3),
        no en el chat del orquestador.
  \item \textbf{Autocuración:} la guardia de mtime evita que un MCP ``recupere''
        contexto basura tras un flujo que no llegó a escribir estado; el
        \texttt{bitacoras.py check} evita cerrar una sesión sin dejar rastro semántico.
  \item \textbf{Retry budget respetado:} sin reintentos infinitos; el criterio es
        binario y verificable.
\end{itemize}

\begin{cajaEjemplo}
\faIcon{code}~\textbf{Ejemplo de uso cotidiano:} tras el reinicio de una sesión (o
antes de reanudar un flujo multi-paso), basta con
\begin{center}
\texttt{python3 execution/estado\_sesion.py check}
\end{center}
para saber si \texttt{run\_state.json} es fiable. Si reporta \texttt{huerfano}, un
\texttt{clean} lo purga en un paso.
\end{cajaEjemplo}

% ══ SECCIÓN 8: PENDIENTES ═════════════════════════════════════════════════════
\section{\iconotexto{list}{Pendientes y líneas siguientes}}

\begin{itemize}[leftmargin=*]
  \item Considerar que los flujos emitan su \texttt{run\_id} al final del stdout
        (permitiría validación robusta sin depender de mtime).
  \item Extender \texttt{check} para limpiar \texttt{run\_state\_*.json} \emph{por
        run} de forma aún más agresiva solo si el usuario lo solicita.
  \item Integrar \texttt{bitacoras.py check} como recordatorio automático de cierre
        de sesión (en flujo de arranque o directiva).
  \item El pendiente de \texttt{docs/AGENTE\_IA/orquestador\_repo\_a\_skill.mp4}
        permanece en manos del usuario (fuera del alcance de esta sesión).
\end{itemize}

\end{document}
